\section{Managerial Process Plan}

Three people building twelve interconnected modules still need explicit rules for coordination, review, and change management \cite{spp1}. Without them, effort duplicates, documentation drifts, and scope expands invisibly.

\subsection{Communication and Coordination}

The team meets at the start and end of each sprint. The start-of-sprint meeting sets goals for that four-week block. The end-of-sprint review confirms what shipped, what did not, and why. Between these two fixed points, daily check-ins happen asynchronously. Any blocker that is not resolved in 24 hours escalates to a team discussion \cite{scrum1}.

Documentation is not separated from implementation. Every module that ships generates a documentation update in the same sprint. This prevents the common FYP problem where the report describes the plan while the prototype has moved somewhere entirely different.

\subsection{Review and Feedback Integration}

Supervisor feedback is the highest-priority input in the project. When the supervisor identifies a gap in the SRS, a design inconsistency, or a scope misalignment, the team addresses it in the next sprint. No feedback is deferred to the final two weeks \cite{agile1}.

Internal review happens at the module level. Before a module moves to integration, the team runs it through a checklist: Does it match its SRS requirement? Does it compile and run cleanly? Does it have at least one passing test case? If any of these answers is no, the module does not move forward.

\subsection{Change Management}

Agile expects change. Not every change is worth accepting \cite{agile1}. The team evaluates proposed adjustments against three criteria: Does it improve the academic value of the project? Does it fit within the remaining sprint budget? Does it risk destabilising a module that is already working?

If a change passes all three checks, it enters the next sprint's backlog. If it fails any criterion, it is logged in a future-work register rather than accepted immediately. This keeps the core deliverables stable while remaining genuinely open to improvement.